% Official solution labels these parts B1--B6 (question 2 = section B of the
% 2012 long-question solutions); they correspond to parts 2.1--2.6 of the
% statement. The A section in the same crop belongs to long question 1 and is
% not transcribed here.
\begin{description}
\item[B1.] \mbox{}
  \begin{align*}
    &1^{\mathrm{st}}\ \text{line:}\quad \{(0.0);\,(15.210)\}\ \
      (10^3\,\mathrm{ly},\ \mathrm{km/s}) \\
    &2^{\mathrm{nd}}\ \text{line:}\quad \{(15.210);\,(60.200)\}\ \
      (10^3\,\mathrm{ly},\ \mathrm{km/s})
  \end{align*}
  \[ V(D) = \begin{cases} 14D, & 0 \le D < 15 \\
       213 - 0.22D, & 15 \le D \le 60 \end{cases}, \qquad
     D \text{ in } 10^3\,\mathrm{l.y.} \text{ and } V \text{ in km/s} \]

\item[B2.] Using the equations from (a):
  \[ \omega(D) = \frac{V(D)}{D} \;\Rightarrow\;
     \omega(D) = \begin{cases} 14, & 0 \le D < 15 \\[2pt]
       \dfrac{213}{D} - 0.22, & 15 \le D \le 60 \end{cases}
     \quad \frac{\mathrm{km/s}}{10^3\,\mathrm{l.y.}} \]
  \[ \omega_{\mathrm{wave}}(D) = 2\omega(D) \]
  The time that a spiral arm takes to have one more turn around the galactic
  center:
  \[ \omega_{\mathrm{rel}} = \omega_{\mathrm{wave}}^{\mathrm{max}}
     - \omega_{\mathrm{wave}}^{\mathrm{min}} \]
  \[ \omega_{\mathrm{rel}} = \frac{1}{2}\left[
       \left(\frac{213}{15} - 0.22\right)
       - \left(\frac{213}{80} - 0.22\right)\right]
     = 23.9\,\frac{\mathrm{km/s}}{10^3\,\mathrm{l.y.}}
     % sic: as printed; the prefactor 1/2 and the denominator 80 do not
     % reproduce 23.9 (1/2 x 11.54 = 5.77), which is closer to 2 x 11.54
  \]
  \[ P_{\mathrm{sin}} = \frac{2\pi}{\omega_{\mathrm{rel}}}
     = 8.15\cdot 10^{7}\ \text{years} \]

\item[B3.] Tully--Fisher relation estimates the luminosity of the galaxy, and
  considering $\Delta V = 2V_{\mathrm{max}}$:
  \[ L = 16\,k\,V_{\mathrm{max}}^{4} \]
  We can compare the magnitude of the galaxy with the Sun:
  \[ m - M = 5\log d - 5 \;\Rightarrow\;
     m - \left(M_{\mathrm{Sun}}
       - 2.5\log\frac{L}{L_{\mathrm{Sun}}}\right) = 5\log d - 5 \]
  \[ \log d = 1 + \frac{m - M_{\mathrm{Sun}}
       + 2.5\log\!\left(\dfrac{16\,k\,V_{\mathrm{max}}^{4}}
       {L_{\mathrm{Sun}}}\right)}{5} \]
  From the graph: $V_{\mathrm{max}} = 230$\,km/s:
  \[ \log d = 1 + \frac{8.5 - 4.8
       + 2.5\log\!\left(16\cdot 0.317\cdot 230^{4}\right)}{5}
     = 6.816 \;\Rightarrow\; d = 6.5\,\mathrm{Mpc} \]

\item[B4.] We can estimate the recession velocity of the galaxy using
  Hubble's law:
  \[ V_0 = H\cdot d \]
  The maximum and minimum wavelengths occur where the radial velocity is
  $V_{\mathrm{max}}$, the radial velocity is:
  \[ V = V_0 \pm V_{\mathrm{max}} \]
  Using $\dfrac{\Delta\lambda}{\lambda} = \dfrac{v}{c}$, the wavelengths are:
  \begin{align*}
    \lambda_{\mathrm{max}} &= \lambda\left(1
      + \frac{Hd + V_{\mathrm{max}}}{c}\right) \\
    \lambda_{\mathrm{min}} &= \lambda\left(1
      + \frac{Hd - V_{\mathrm{max}}}{c}\right)
  \end{align*}
  Calculating the values, using $\lambda = 656.28$\,nm:
  \begin{align*}
    \lambda_{\mathrm{max}} &= 657.79\,\mathrm{nm} \\
    \lambda_{\mathrm{min}} &= 656.78\,\mathrm{nm}
  \end{align*}

\item[B5.] Supposing a mass with spherical symmetry distribution and a
  circular orbit, we have:
  \[ \frac{V^2}{D} = \frac{G M_{\mathrm{int}}}{D^2} \;\Rightarrow\;
     M_{\mathrm{int}} = \frac{V^2 D}{G} \]
  From the graph $V(3.0\cdot 10^{4}\ \mathrm{l.y.}) = 225$\,km/s, So the mass
  is:
  \[ M_{\mathrm{int}} = 2.15\cdot 10^{41}\,\mathrm{kg}
     = 1.08\cdot 10^{11}\,M_{\mathrm{sun}} \]

\item[B6.] As the mass is derived from the rotation curve, it is necessary to
  take into account that part of the mass is dark matter, and only part of the
  mass is in the stars ($\rho_{\mathrm{star}} = 1/3$). The fraction of
  baryonic matter is:
  \[ \rho_{\mathrm{matter}} = \frac{f_{\mathrm{baryon}}}{f_{\mathrm{matter}}}
     = \frac{4}{22+4} = \frac{4}{26} \]
  The number of stars is then:
  \[ N = \rho_{\mathrm{matter}}\cdot\rho_{\mathrm{star}}\,
     \frac{M_{\mathrm{galaxy}}}{M_{\mathrm{Sun}}}
     = \frac{4}{26}\,\frac{1}{3}\,1.08\cdot 10^{11}
     = 5.5\cdot 10^{9}\ \text{stars} \]
\end{description}
